% 第 16 题(2)：防窥屏结构与可视角度（TikZ 线稿）
\begin{tikzpicture}[font=\small, >=Stealth, scale=0.82]
  % 透明介质
  \fill[gray!12] (0,1.5) rectangle (8.2,2.8);
  \draw[thick] (0,1.5) -- (8.2,1.5) -- (8.2,2.8) -- (0,2.8) -- cycle;
  % 屏障
  \foreach \x in {2.0,3.7,5.4,7.1}{
    \draw[line width=2.6pt] (\x,1.5) -- (\x,2.8);
  }
  % 发光像素单元
  \draw[line width=3pt] (4.35,1.5) -- (4.75,1.5);
  % 像素发出的光线（到两侧屏障顶端）
  \draw[->, thick] (4.55,1.55) -- (3.7,2.8);
  \draw[->, thick] (4.55,1.55) -- (5.4,2.8);
  % 折射到空气中的光线
  \draw[->, thick] (3.7,2.8) -- ++(150:1.0);
  \draw[->, thick] (5.4,2.8) -- ++(30:1.0);
  \draw[dashed] (3.7,2.8) -- (3.7,3.7);
  \draw[dashed] (5.4,2.8) -- (5.4,3.7);
  % 可视角度 θ（两条折射光线之间的夹角）
  \draw (4.55,3.0) arc[start angle=145, end angle=35, radius=0.95];
  \node at (4.55,3.62) {$\theta$};
  % 标注
  \node at (0.95,3.02) {空气};
  \node at (0.95,2.12) {透明介质};
  \draw[<->] (3.7,1.30) -- (5.4,1.30) node[midway, below] {$l$};
  \draw[<->] (7.65,1.5) -- (7.65,2.8) node[midway, right] {$d$};
  \node[right] at (5.70,1.02) {发光像素单元};
  \draw[->] (5.65,1.08) -- (4.82,1.48);
\end{tikzpicture}
